\section*{1. Wstęp historyczny}
Badany obiekt pochodzi z okresu późnego baroku i wykazuje liczne zmiany autorskie oraz wtórne przemalowania.

\section*{2. Wyniki analizy pigmentów}
\begin{tabularx}{\textwidth}{lXX}
\toprule
\textbf{Warstwa} & \textbf{Zidentyfikowany pigment} & \textbf{Stan zachowania} \\
\midrule
Podmalowanie & Biel ołowiana z węglem drzewnym & Stabilny, spękania siatkowe \\
Warstwa lica & Ultramaryna naturalna i cynober & Miejscowe blaknięcia i korozja spoiwa \\
Werniks & Żywica damarowa (wtórna) & Zżółknięty, wymaga usunięcia \\
\bottomrule
\end{tabularx}

\section*{3. Zalecenia konserwatorskie}
Zaleca się bezpieczne usunięcie zżółknionego werniksu metodami chemicznymi oraz punktowe uzupełnienie ubytków spoiwa.

\vspace{1cm}
\noindent\textbf{Dokumentacja:} Tobiasz Lorentz